\documentclass{article}

\usepackage{amsmath}
\usepackage{booktabs}
\usepackage{enumitem}
\usepackage[most]{tcolorbox}

\newtcolorbox{eqbox}{colback=blue!6!white, colframe=blue!45!black,
  boxrule=0.8pt, sharp corners, left=2pt, right=2pt, top=2pt, bottom=2pt}
\newtcolorbox{whybox}{colback=black!4!white, colframe=black!4!white,
  borderline west={2pt}{0pt}{black!30}, sharp corners,
  left=6pt, right=6pt, top=4pt, bottom=4pt}

% Compact variable-definition list placed under each equation
\newlist{vardefs}{itemize}{1}
\setlist[vardefs]{label=--, nosep, leftmargin=1.6em, topsep=2pt}

\title{Dynamic No-Trade Bands for Intraday Delta Hedging}
\author{Alex Marthan}
\date{July 2026}

\begin{document}
\maketitle

\section{Model Framework}

\subsection{From a Static Band to a State-Dependent Band}

The static strategy holds a constant no-trade band $H$ around
$\Delta_{\text{target}} = 0$ and hedges only when $\lvert\Delta(t)\rvert$ clears it. A
constant $H$ encodes a constant answer to the question the band actually poses: \emph{how
much residual delta is worth carrying rather than paying to remove?} That answer is not
constant. The cost of carrying delta over the next interval depends on how fast the book
regenerates delta once hedged, and that speed is a measurable state variable.

The dynamic model therefore replaces the constant with a per-bar band:

\begin{eqbox}
\begin{equation}
-H(t) \le \Delta(t) \le H(t),
\qquad
\lvert \Delta(t) \rvert > H(t)\,(1+\epsilon)
\ \Rightarrow\ \text{hedge}.
\label{eq:dyntrigger}
\end{equation}
\end{eqbox}

Everything else in the strategy --- execution scheduling, cooldown, cost and risk
measurement, benchmarks --- is unchanged. The static rule is the special case
$H(t) \equiv H_0$, so the dynamic model strictly nests it and any improvement must be
earned.

\subsection{The Gamma-Conditioned Band}

Between trade arrivals the book's delta does not sit still: it drifts with the underlying
at a rate set by the book's convexity, $d\Delta \approx \Gamma\,dS$. Gamma is thus the
\emph{speed} at which market moves regenerate delta after a hedge. For a band of
half-width $H$ and per-bar underlying moves of scale $\sigma_S$, the expected number of
bars for gamma drift alone to carry delta across the band is of order

\[
\left( \frac{H}{\lvert\Gamma\rvert\,\sigma_S} \right)^{2}.
\]

A high-gamma book traverses a fixed band quickly and therefore carries more unhedged risk
between checks; a low-gamma book barely moves and the same band is needlessly tight.
Making the band shrink as gamma rises equalizes this breach behaviour across regimes:

\begin{eqbox}
\begin{equation}
H(t) = \frac{H_0}{1 + c\,\lvert \Gamma(t) \rvert}.
\label{eq:gammaband}
\end{equation}
\end{eqbox}

\begin{vardefs}
\item $H_0$: base band half-width, in hedge-instrument units --- the band the rule uses
  when the book is gamma-flat;
\item $\Gamma(t)$: book gamma at bar $t$, in hedge-instrument units per unit move of the
  underlying;
\item $c$: tilt strength, controlling how hard the band responds to gamma.
\end{vardefs}

\textbf{Units of $c$.} Gamma enters undivided, so $c$ is \emph{not} dimensionless: it
carries the reciprocal units of $\Gamma$, and the band halves exactly when
$\lvert\Gamma\rvert = 1/c$. The useful range of $c$ is therefore set by the book's own
gamma scale and differs by orders of magnitude across instruments
(Table~\ref{tab:instrument}). Dividing $\lvert\Gamma\rvert$ by its sample median and
rescaling $c$ gives the identical family of bands with a dimensionless $c$; only the
units of the parameter change, not the model.

\textbf{Limiting cases.} $c \to 0$ recovers the static band $H(t) \equiv H_0$; large $c$
drives $H(t) \to 0$ whenever the book carries gamma, approaching continuous rehedging.
The backtest selects $c$, so the data decides how hard to tilt.

\begin{whybox}
\textbf{The sign of the effect is not obvious.} Classical optimal-hedging asymptotics
under proportional transaction costs (Whalley--Wilmott; Zakamouline) give an optimal
half-width $\propto \Gamma^{2/3}$ --- i.e.\ \emph{wider} when gamma is high, on the
grounds that chasing a fast-moving delta churns execution costs. Equation~\eqref{eq:gammaband}
tilts the other way, prioritising risk control over turnover. Which effect dominates
depends on the cost model and the risk measure, so $c$ is gridded and $c \to 0$ recovers
the static band: the comparison is decided empirically, not assumed.
\end{whybox}

\subsection{Book Gamma Revaluation}

The band in Equation~\eqref{eq:gammaband} requires $\Gamma(t)$ at every bar. It is built
exactly like the book delta, by repricing every live position at the current underlying
price and aggregating:

\begin{eqbox}
\begin{equation}
\Gamma(t_j) = \sum_{i:\,t_i \le t_j} s_i \, n_i \, \gamma_i(t_j) \, m,
\qquad
\gamma_i(t_j) = \frac{\varphi\!\left(d_{1,i}(t_j)\right)}
                     {S_j\,\sigma\,\sqrt{\tau_i(t_j)}}.
\label{eq:gamma}
\end{equation}
\end{eqbox}

\begin{vardefs}
\item $\gamma_i(t_j)$: per-contract gamma of option $i$ at bar $t_j$ (identical for calls
  and puts);
\item $\varphi(\cdot)$: standard normal density; $d_{1,i}$ is as in the delta definition;
\item $S_j$: underlying price at bar $t_j$ (futures price $F_j$ for ES);
\item $m$: contract-to-hedge-unit multiplier ($m = 1$ for ES options, whose deltas are
  already in futures equivalents; $m = 100$ for equity options, which deliver 100 shares);
\item $s_i, n_i, \sigma, \tau_i$: as in the delta definition --- trade sign, contract
  count, the day's flat volatility input, and time to expiry.
\end{vardefs}

\textbf{Gamma is strategy-independent.} The hedge instrument --- a future or the stock
itself --- has zero gamma, so hedging changes $\Delta$ but never $\Gamma$. The path
$\Gamma(t)$ therefore depends only on the option tape and the price path, not on the
hedging rule, and is precomputed once per day and reused by every configuration in the
grid. This is what makes the dynamic grid search no more expensive than the static one.

\subsection{Instrument Instantiation}

The framework above is instrument-agnostic; only the units and the gamma scale change.

\begin{table}[htbp]
\centering
\small
\begin{tabular}{lll}
\toprule
 & \textbf{ES futures} & \textbf{SPCX equity} \\
\midrule
Hedge instrument      & ES future            & The stock \\
Delta / band unit     & Contract equivalents & Shares \\
Contract multiplier $m$ & $1$                & $100$ \\
P\&L multiplier $M_h$ & $50$ \$/point        & $1$ \$/share \\
Option delta model    & Black-76 (futures)   & Black--Scholes (spot) \\
$\Gamma$ units        & contracts/point      & shares/\$ \\
Median $\lvert\Gamma\rvert$ & $0.0928$       & $2{,}176$ \\
$c$ units             & points/contract      & \$/share \\
$c$ grid              & $\{5, 10, 20\}$      & $\{2,5,10,20\}\times10^{-4}$ \\
\bottomrule
\end{tabular}
\caption{Instantiation of the dynamic band model on the two books. The $c$ grids are
\emph{not} interchangeable: each is scaled to its own book's gamma so that
$c\,\lvert\Gamma\rvert$ spans roughly $0.5$--$2$.}
\label{tab:instrument}
\end{table}

\section{Strategy and Experimental Design}

This section states the experimental protocol shared by both books. Only the parameter
\emph{values} in Table~\ref{tab:grid} differ between them.

\subsection{Trigger Rule and Hedge Sizing}

The band is monitored on a 5-minute grid. A hedge fires when
Equation~\eqref{eq:dyntrigger} is satisfied, and the parent order returns the book to
flat:

\begin{eqbox}
\begin{equation}
\begin{aligned}
Q_t &= -\Delta(t) && \text{(full flatten)},\\[2pt]
Q_t &= -\,\mathrm{sign}(\Delta(t))\,\big(\lvert\Delta(t)\rvert - H(t)\big)
     && \text{(edge target)}.
\end{aligned}
\label{eq:sizing}
\end{equation}
\end{eqbox}

The edge-target variant trades only the excess over the band and leaves the residual
parked at the edge; it is evaluated as a refinement in both results sections. After a
schedule completes the trigger is suppressed for $T_{\text{cooldown}}$ unless the breach
is large ($\lvert\Delta(t)\rvert > 2H(t)$, with the override threshold scaling with the
\emph{current} band). Any residual delta is forcibly flattened at the close, so every
configuration satisfies $\Delta(T) = 0$.

\subsection{Execution}

The parent order is split into child orders over a horizon $\tau$, i.e.\ over
$N = \tau/5\,\text{min}$ consecutive bars, either equally weighted (TWAP) or weighted by
each bar's share of traded volume (VWAP). A new trigger cancels and replaces any running
schedule, resized to the current residual delta.

\subsection{Cost and Risk Measurement}

Execution cost is a market-data proxy combining a half-spread term and a concave impact
term:

\begin{eqbox}
\begin{equation}
C_{\text{exec}} = \sum_k M_h \lvert q_k \rvert
\left[ \frac{s}{2} + \alpha\,S_k \left( \frac{\lvert q_k \rvert}{L_k} \right)^{\!\beta} \right],
\qquad \beta = 0.5 .
\label{eq:cost}
\end{equation}
\end{eqbox}

\begin{vardefs}
\item $q_k$: signed child order in bar $k$; $S_k$ the underlying price there;
\item $s/2$: half bid--ask spread in price units;
\item $\alpha$: impact level; $L_k$ the liquidity the fill leans on.
\end{vardefs}

\textbf{$\alpha$ and $L_k$ are the one place the two books genuinely diverge}, and the
divergence is deliberate rather than a porting artifact. On ES, $\alpha$ is anchored so a
1-lot pays $\approx0.5$bp of notional at median ADV, and $L_k$ is full-day ADV --- both
harmless for an instrument whose hedge clips are $\sim\!10^{-4}$ of ADV. Neither survives
on SPCX, where clips run to tens of percent of ADV: there $\alpha = Y\sigma_{\text{daily}}$
follows the square-root law of impact (Toth--Bouchaud, $Y = 0.5$), anchoring impact to the
book's own volatility, and $L_k$ is the liquidity of the fill's \emph{execution window}
($\text{ADV}\times 5/390$ for an intraday child) rather than the whole day. Both books
treat the final bar as an auction cross: zero spread, full-day depth, and an impact
haircut.

Delta risk is the root mean square of realized residual-delta P\&L per bar:

\begin{eqbox}
\begin{equation}
\ell^{\Delta}_{j+1} = M_h\,\Delta(t_j)\,(S_{j+1} - S_j),
\qquad
C_{\text{risk}} = \sqrt{ \frac{1}{M} \sum_{j=0}^{M-1} \big(\ell^{\Delta}_{j+1}\big)^{2} }.
\end{equation}
\end{eqbox}

\subsection{Benchmarks and Standardized Objective}

\textbf{Zero Hedge} carries the book all day and flattens once into the close (cost floor,
risk ceiling); \textbf{Full Hedge} re-flattens every bar (risk floor, cost ceiling). Both
bound the attainable region and normalize the objective:

\begin{eqbox}
\begin{equation}
L_{\text{std}}(\theta) =
\frac{C_{\text{exec}}(\theta)}{C_{\text{exec}}(\text{Full Hedge})}
+ \lambda_{\text{risk}}\,
\frac{C_{\text{risk}}(\theta)}{C_{\text{risk}}(\text{Zero Hedge})},
\qquad \lambda_{\text{risk}} = 1 .
\end{equation}
\end{eqbox}

\subsection{Fair-Comparison Methodology}

Two protocol choices carry most of the inferential weight, because the naive comparison
is misleading.

\textbf{Width matching.} At the same $H_0$ the gamma band is \emph{narrower on average}
than the static band, since the divisor in Equation~\eqref{eq:gammaband} is $\ge 1$. A
comparison at matched $H_0$ therefore conflates two different things: the value of
\emph{conditioning the band on gamma} and the value of simply running a \emph{tighter}
band. Every headline dynamic-vs-static claim below is instead made against a static band
set to the gamma configuration's \emph{mean effective width}, under identical
$\epsilon$, cooldown, horizon and method.

\textbf{Hedge-frequency capping.} The cost proxy in Equation~\eqref{eq:cost} contains no
fixed per-ticket charge, and impact is concave in clip size, so splitting a hedge into
more, smaller child orders is strictly cheaper in the model with no floor. Unconstrained
grid search consequently runs away toward the tightest band available --- a boundary
solution, not a strategy. Capping the average number of hedges per day stands in for the
missing per-ticket cost and restores an interior, decision-relevant optimum.

\subsection{Parameter Grid}

\begin{table}[htbp]
\centering
\footnotesize
\setlength{\tabcolsep}{4pt}
\begin{tabular}{llll}
\toprule
\textbf{Parameter} & \textbf{Role} & \textbf{ES grid} & \textbf{SPCX grid} \\
\midrule
$H_0$ & Band width & $\{2,5,10,25,50,100,200\}$ ct & $\{5,10,25,50,100,200,300\}$k sh \\
$c$   & Gamma tilt & $\{5, 10, 20\}$               & $\{2,5,10,20\}\times10^{-4}$ \\
$\tau$ & Horizon   & \multicolumn{2}{l}{$\{15, 30, 60\}$ min} \\
$e$   & Method     & \multicolumn{2}{l}{$\{\text{TWAP}, \text{VWAP}\}$} \\
$\epsilon$ & Buffer & \multicolumn{2}{l}{$5\%$ (held fixed)} \\
$T_{\text{cooldown}}$ & Cooldown & \multicolumn{2}{l}{$15$ min (held fixed)} \\
\midrule
$\lambda_{\text{risk}}$ & Risk weight & \multicolumn{2}{l}{$\{0.25, 0.5, 1, 2, 4\}$ (sweep)} \\
\bottomrule
\end{tabular}
\caption{Strategy grid. $\epsilon$ and $T_{\text{cooldown}}$ are held fixed across all
configurations so that the static-versus-dynamic comparison is not confounded by the
second-order trigger refinements.}
\label{tab:grid}
\end{table}

Each configuration is simulated over every day in the sample using the same precomputed
delta and gamma paths, so differences across configurations reflect the hedging rule
alone.

\section{Results: ES Futures}

\textbf{Data.} 102 trading days (January 2 -- May 29, 2026), a simulated tape of
32{,}127 executed ES option trades, and 8{,}364 five-minute futures bars. Cost model:
$\alpha = 0.055$, half-spread $0.125$ points, closing-auction impact factor $0.15$,
$M_h = 50$.

\begin{table}[htbp]
\centering
\small
\begin{tabular}{lrrr}
\toprule
\textbf{Benchmark} & \textbf{Exec.\ cost} & \textbf{Risk (RMS/bar)} & $L_{\text{std}}$ \\
\midrule
Zero Hedge & \$235{,}240   & \$24{,}801 & 1.058 \\
Full Hedge & \$4{,}030{,}988 & \$0      & 1.000 \\
\bottomrule
\end{tabular}
\caption{ES benchmark bounds over 102 days.}
\end{table}

\textbf{Gamma signal.} Book gamma has a median of $0.0928$ contracts/point (90th/99th
percentiles $0.368$ / $0.999$). Within a day $\lvert\Gamma\rvert$ has a coefficient of
variation of $0.61$ and a max/median ratio of $2.1\times$; across bars it correlates with
$\lvert\Delta\rvert$ at $\rho = 0.63$. Gamma is thus a strongly varying but only
partially independent signal on this book.

\begin{table}[htbp]
\centering
\small
\begin{tabular}{lrrrrr}
\toprule
\textbf{Family} & $H_0$ & $c$ & \textbf{Exec.\ cost} & \textbf{Risk} & $L_{\text{std}}$ \\
\midrule
Static      & 2  & ---  & \$974{,}803 & \$6{,}270 & 0.4946 \\
Combined    & 10 & 20   & \$653{,}688 & \$8{,}309 & 0.4972 \\
Gamma       & 10 & 20   & \$668{,}157 & \$8{,}320 & 0.5012 \\
Time-decay  & 10 & ---  & \$721{,}897 & \$8{,}331 & 0.5150 \\
\bottomrule
\end{tabular}
\caption{Best configuration per band family (TWAP, 30--60min). ``Combined'' multiplies the
gamma band by a $\sqrt{R(t)/R(0)}$ time-decay factor; ``time-decay'' uses that factor
alone.}
\end{table}

\begin{whybox}
\textbf{The width trap.} The unconstrained winner is the \emph{tightest static band in the
grid} ($H_0 = 2$, $L_{\text{std}} = 0.4946$), and the static ladder falls monotonically
into that grid edge ($H_0 = 25$: $0.5568$; $10$: $0.5321$; $5$: $0.5068$; $2$: $0.4947$).
This is the boundary solution of Section~2.6, not evidence that dynamics fail: at matched
$H_0$ the dynamic families look far better than static purely because they are narrower on
average. Only the width-matched comparison separates the two effects.
\end{whybox}

\textbf{Width-matched comparison.} Setting a static band to each dynamic configuration's
mean effective width isolates the timing information, and it splits the two mechanisms
cleanly.

\begin{table}[htbp]
\centering
\small
\begin{tabular}{rrrrr}
\toprule
$H_0$ & \textbf{Mean width} & \textbf{Gamma} & \textbf{Combined} & \textbf{Time-decay} \\
\midrule
10  & $4.1$  & $+0.22\%$ & $+0.34\%$ & $-0.22\%$ \\
25  & $10.3$ & $+2.56\%$ & $+1.44\%$ & $-0.25\%$ \\
50  & $20.6$ & $+3.23\%$ & $+4.38\%$ & $-1.35\%$ \\
100 & $41.1$ & $+4.02\%$ & $+0.83\%$ & $-8.52\%$ \\
200 & $82.3$ & $+2.31\%$ & $+2.32\%$ & $-3.31\%$ \\
\bottomrule
\end{tabular}
\caption{Edge over a width-matched static band, \% of $L_{\text{std}}$ (TWAP 60min).
Positive = dynamic better. Mean width in contract equivalents.}
\end{table}

Gamma conditioning holds a genuine edge of roughly $2.5$--$4.0\%$ at moderate widths. The
time-decay factor is \emph{negative} at every width: its rationale was that the forced
end-of-day flatten is expensive, but the closing-auction discount removes exactly that
penalty, so shrinking the band into the close pre-pays execution for nothing. The combined
form inherits a diluted positive edge from its gamma component.

\textbf{Robustness.} On three re-seeded tapes (per-day bootstrap re-drawing contract, side
and size, with $c$ and the matched widths frozen), the gamma edge is positive on every
seed ($+1.18\%$, $+0.58\%$, $+0.47\%$), as is the combined form; time-decay remains mixed.

\textbf{Band-edge targeting.} Replacing the full flatten with the edge target of
Equation~\eqref{eq:sizing} is the strongest single refinement on ES.

\begin{table}[htbp]
\centering
\small
\begin{tabular}{lrrrrr}
\toprule
$H_0$ & 5 & 10 & 25 & 50 & 100 \\
\midrule
Static band      & $+2.6\%$ & $+5.4\%$ & $+6.0\%$ & $+1.5\%$ & $-9.4\%$ \\
Combined dynamic & $+0.3\%$ & $+0.8\%$ & $+1.8\%$ & $+1.7\%$ & $+3.7\%$ \\
\bottomrule
\end{tabular}
\caption{Improvement from edge targeting (TWAP 60min); positive = better than full
flatten.}
\end{table}

Parking residual delta at the band edge and carrying it to a cheap auction close is worth
up to $+6\%$ for moderate static widths, though it backfires for wide static bands, where
too much risk is parked under $\lambda_{\text{risk}} = 1$. The dynamic band removes that
failure mode --- its shrinking width pulls the parked delta inward --- and the best overall
configuration is the \textbf{combined dynamic band with edge targeting}
($H_0 = 5$, $c = 20$: $L_{\text{std}} = 0.4907$, against $0.4946$ for the best
full-flatten configuration).

\section{Results: SPCX Equity}

\textbf{Data.} 22 trading days (June 16 -- July 17, 2026), 120{,}951 simulated SPCX option
executions, 1{,}716 five-minute bars. The stock traded \$123.12--\$222.80 (median
\$155.27) with ATM implied volatility of 70--126\% and a 20-day ADV proxy near
982{,}000 shares. Cost model: measured half-spread \$0.027/share,
$\alpha = Y\sigma_{\text{daily}} = 0.0273$, execution-window liquidity
$L_k = \text{ADV}\times 5/390$, closing-cross impact factor $0.30$, $M_h = 1$.

\begin{table}[htbp]
\centering
\small
\begin{tabular}{lrrr}
\toprule
\textbf{Benchmark} & \textbf{Exec.\ cost} & \textbf{Risk (RMS/bar)} & $L_{\text{std}}$ \\
\midrule
Zero Hedge & \$9{,}974{,}886   & \$435{,}115 & 1.027 \\
Full Hedge & \$367{,}974{,}631 & \$0         & 1.000 \\
\bottomrule
\end{tabular}
\caption{SPCX benchmark bounds over 22 days. The selected strategy turns over
$\approx638{,}000$ shares/day, about 65\% of ADV, so impact rather than spread dominates.}
\end{table}

\textbf{Only the gamma mechanism is carried over.} The time-decay factor is dropped for
this book: the Nasdaq closing cross makes the end-of-day flatten cheap, removing the
penalty that factor was designed to avoid --- the same effect that turned its ES edge
negative.

\textbf{Gamma signal.} Median $\lvert\Gamma\rvert$ is $2{,}176$ shares/\$, with a
within-day coefficient of variation of $0.46$ and a max/median ratio of $2.1\times$.
Across bars gamma correlates with $\lvert\Delta\rvert$ at only $\rho = 0.21$ --- more
independent information than on ES, but a narrower range over which the band can flex.

\begin{table}[htbp]
\centering
\small
\begin{tabular}{lrrrrr}
\toprule
\textbf{Family} & $H_0$ (sh) & $c$ & \textbf{Exec.\ cost} & \textbf{Risk} & $L_{\text{std}}$ \\
\midrule
Static & 5{,}000  & ---                & \$88.5M & \$112{,}916 & 0.5001 \\
Gamma  & 25{,}000 & $2\times10^{-3}$   & \$88.2M & \$113{,}396 & 0.5004 \\
\bottomrule
\end{tabular}
\caption{Best configuration per family (VWAP, 30min). Both cut the objective by
$\approx51\%$ against Zero Hedge and $\approx50\%$ against Full Hedge --- the first-order
result on this book.}
\end{table}

\begin{whybox}
\textbf{The same width trap, and the same boundary solution.} The unconstrained winner is
again the tightest static band in the grid, firing $47.6$ hedges/day, with the static
ladder falling monotonically into the edge ($50$k: $0.5287$; $25$k: $0.5238$; $10$k:
$0.5088$; $5$k: $0.5001$). No desk trading a book that turns over 65\% of ADV would hedge
that often, so the unconstrained optimum is not a usable comparison point.
\end{whybox}

\textbf{Width-matched comparison.} Width-for-width the gamma band is risk-efficient across
every band width that actually binds, and the edge \emph{grows} with width: wider bands
bind less often, so the choice of \emph{when} to hedge matters more.

\begin{table}[htbp]
\centering
\small
\begin{tabular}{rrrrr}
\toprule
$H_0$ (sh) & \textbf{Mean width} & \textbf{Gamma} $L$ & \textbf{Matched static} $L$ & \textbf{Edge} \\
\midrule
25{,}000  & 9{,}049   & 0.5243 & 0.5248 & $+0.11\%$ \\
50{,}000  & 18{,}098  & 0.5268 & 0.5324 & $+1.06\%$ \\
100{,}000 & 36{,}195  & 0.5288 & 0.5386 & $+1.82\%$ \\
200{,}000 & 72{,}391  & 0.5292 & 0.5480 & $+3.42\%$ \\
300{,}000 & 108{,}586 & 0.5904 & 0.5389 & $-9.55\%$ \\
\bottomrule
\end{tabular}
\caption{Width-matched gamma vs.\ static at $c = 10^{-3}$, TWAP 60min. The edge collapses
at the widest setting, where the band rarely binds at all and the tilt merely
over-tightens on the few occasions it does.}
\end{table}

\textbf{The operational regime.} Capping the hedge budget --- the stand-in for the missing
per-ticket cost --- is what makes the comparison decision-relevant, and it is where the
gamma band wins outright.

\begin{table}[htbp]
\centering
\small
\begin{tabular}{lrrrl}
\toprule
\textbf{Cap/day} & \textbf{Static} $L$ & \textbf{Gamma} $L$ & \textbf{Edge} & \textbf{Winner} \\
\midrule
5    & 0.5242 & 0.5444 & $-3.86\%$ & static \\
10   & 0.5242 & 0.5338 & $-1.83\%$ & static \\
15   & 0.5242 & 0.5241 & $+0.01\%$ & level \\
20   & 0.5238 & 0.5118 & $+2.29\%$ & gamma \\
30   & 0.5238 & 0.5086 & $+2.91\%$ & gamma \\
\midrule
none & 0.5001 & 0.5004 & $-0.06\%$ & static \\
\bottomrule
\end{tabular}
\caption{Best static vs.\ best gamma under a cap on hedge frequency.}
\end{table}

Below $\approx10$ hedges/day both families are forced so wide that the band rarely binds
and shape barely matters, so the simpler static rule wins. The two draw level at
$\approx15$/day, and from $\approx20$/day upward --- the realistic operating range --- the
gamma band wins by $+2.3\%$ to $+2.9\%$. A fixed hedge budget is precisely the constraint
that makes \emph{where} hedges are spent matter, which is what
Equation~\eqref{eq:gammaband} optimizes. The advantage is a capacity-regime effect, not an
artifact of the risk weight: unconstrained, the degenerate tight static band wins at every
$\lambda_{\text{risk}} \in [0.25, 4]$.

\textbf{Band-edge targeting.} Unlike on ES, edge targeting adds little here.

\begin{table}[htbp]
\centering
\small
\begin{tabular}{lrrrrr}
\toprule
$H_0$ (sh) & 10{,}000 & 25{,}000 & 50{,}000 & 100{,}000 & 200{,}000 \\
\midrule
Static band & $+0.59\%$ & $+1.14\%$ & $+2.00\%$ & $+0.31\%$ & $-3.22\%$ \\
Gamma band  & $+0.10\%$ & $+0.36\%$ & $-0.04\%$ & $-0.20\%$ & $-2.43\%$ \\
\bottomrule
\end{tabular}
\caption{Improvement from edge targeting (TWAP 60min).}
\end{table}

Carrying residual delta to a cheap auction close is most of what edge targeting buys, and
on SPCX the closing cross already delivers that to every strategy. Stacked on the gamma
band it adds essentially nothing, turning negative beyond $25{,}000$ shares: the gamma
band already tightens where the residual would otherwise be parked.

\textbf{Comparison with ES.} The two books agree on the substance and differ in degree.
Both show the same boundary solution under a frictionless cost proxy, and on both the
gamma-conditioned band is risk-efficient width-for-width --- $2.5$--$4.0\%$ on ES,
up to $3.4\%$ on SPCX. They differ in what the refinement is worth operationally: on ES,
edge targeting is the larger effect and stacks with the band, whereas on SPCX the closing
cross absorbs that gain and the gamma tilt only becomes decisive once the hedge budget is
capped. The dominant caveats are the absence of a fixed per-ticket cost in both cost
models, and SPCX's short 22-day sample, which is thin for distinguishing effects of a few
percent.

\end{document}
